\section{Yêu cầu đề tài}

\subsection{Mục tiêu thiết kế}
Thiết kế bộ điều khiển I2C Master tuân thủ chuẩn I2C dựa trên kiến trúc RD1139 của Lattice Semiconductor.

Hỗ trợ tốc độ 100 kHz (Standard Mode) và 400 kHz (Fast Mode).

Hỗ trợ các chế độ:
\begin{itemize}[leftmargin=1.2cm]
    \item Ghi dữ liệu (Write)
    \item Đọc dữ liệu (Read)
    \item Địa chỉ slave 7-bit
    \item Tạo và xử lý đầy đủ START, STOP, ACK/NACK, repeated START.
\end{itemize}

Đảm bảo hoạt động ổn định khi tích hợp vào SoC hoặc FPGA.

\subsection{Phạm vi và cấu hình tốc độ}
Trọng tâm báo cáo là kiến trúc \emph{đơn master}; đa master chỉ được nhắc ở mức bus FSM (arbitration, phát hiện START/STOP do master khác).

Tốc độ SCL được chọn thông qua \texttt{i\_clk\_div\_lsb} và \texttt{i\_mode\_reg[2:0]} cùng bộ chia clock trong khối Clock Generator.

\texttt{Processor Interface} ánh xạ thanh ghi điều khiển/trạng thái và tín hiệu ngắt.

\subsection{Ràng buộc thiết kế}
Hoạt động theo chuẩn I2C:
\begin{itemize}[leftmargin=1.2cm]
    \item Tốc độ: 100 kHz (Standard), 400 kHz (Fast Mode).
    \item SDA phải thay đổi khi SCL = LOW.
    \item START: SDA đi từ HIGH $\rightarrow$ LOW khi SCL = HIGH.
    \item STOP: SDA đi từ LOW $\rightarrow$ HIGH khi SCL = HIGH.
    \item Tín hiệu SDA và SCL là open-drain (Inout).
    \item Hoạt động ổn định ở xung clock hệ thống 32 MHz.
    \item Được mô tả bằng Verilog/SystemVerilog HDL với testbench có sẵn.
\end{itemize}

SDA/SCL là \emph{open-drain} với pull-up ngoài; master chỉ kéo mức thấp hoặc để high-Z.

\textit{Clock stretching:} slave có thể giữ SCL thấp — master phải chờ SCL trở lại mức cao trước bit kế tiếp.

\subsection{Giao dịch I2C (tóm tắt)}
Khi bus rảnh, SCL và SDA đều cao nhờ pull-up. Master phát \textbf{START} để chiếm bus. Byte đầu gồm 7 bit địa chỉ và 1 bit R/W (0: master ghi, 1: master đọc). Bit thứ 9: slave kéo SDA thấp là \textbf{ACK}; không kéo là \textbf{NACK} tùy ngữ cảnh. Các byte sau mang dữ liệu hoặc địa chỉ thanh ghi nội bộ của thiết bị slave. Kết thúc bằng \textbf{STOP} hoặc \textbf{Repeated START (Sr)} để mở giao dịch mới mà không nhả bus bằng STOP.

\textbf{Quy tắc timing cốt lõi:} dữ liệu trên SDA chỉ được phép thay đổi khi SCL = LOW; trong pha SCL = HIGH, SDA phải ổn định để phía thu lấy mẫu. Điều kiện START/STOP là ngoại lệ duy nhất được phép thay đổi SDA khi SCL = HIGH.

\textit{\small Tham chiếu chuẩn bus: NXP Semiconductors, \emph{UM10204 I\textsuperscript{2}C-bus specification and user manual}.}

\textbf{Open-drain và mức logic:} SDA/SCL là dây open-drain nên ``mức 1'' đến từ điện trở pull-up; master/slave chỉ có quyền kéo xuống 0 hoặc nhả high-Z. Vì vậy, một số cơ chế như \emph{clock stretching} (slave giữ SCL thấp) và \emph{arbitration} (đa master) dựa trên việc so sánh mức bus thực tế với mức master muốn phát.

\textbf{Các ngoại lệ cần xử lý trong thiết kế:}
\begin{itemize}[leftmargin=1.2cm]
    \item \textbf{NACK}: xảy ra ở địa chỉ hoặc dữ liệu (slave không ACK), cần kết thúc giao dịch và báo lỗi.
    \item \textbf{Clock stretching}: slave giữ SCL thấp $\Rightarrow$ master phải chờ SCL thật về 1 trước khi đi tiếp bit kế.
    \item \textbf{Bus contention/arbitration}: nếu đa master, master phát '1' nhưng bus bị kéo '0' $\Rightarrow$ mất arbitration.
    \item \textbf{Timeout}: SCL bị giữ thấp quá lâu hoặc bus không tiến triển $\Rightarrow$ abort để tránh treo hệ thống.
\end{itemize}

\subsection{Chế độ clock}
\begin{table}[H]
    \centering
    \caption{Chế độ clock và gợi ý thiết kế}
    \label{tab:clock_mode}
    \begin{tabularx}{\textwidth}{|l|l|X|}
        \hline
        \textbf{Chế độ} & \textbf{Tần số SCL (hướng đặt)} & \textbf{Gợi ý} \\
        \hline
        Standard & $\sim$100 kHz & Chu kỳ SCL $\sim$10\,$\mu$s — bộ chia từ 32 MHz cấu hình qua \texttt{i\_clk\_div\_lsb} và \texttt{i\_mode\_reg}. \\
        \hline
        Fast & $\sim$400 kHz & Chu kỳ $\sim$2{,}5\,$\mu$s — margin timing chặt hơn trên pad và trong FSM bit. \\
        \hline
    \end{tabularx}
\end{table}

Với tần số hệ thống $f_{\mathrm{sys}}$ cố định, tần số SCL gần đúng $f_{\mathrm{SCL}} \approx f_{\mathrm{sys}} / (2 \times \mathrm{prescale})$ sau các tầng chia trong Clock Generator.
